\subsection{Union-Exclusion Principle}

The Union-Exclusion Principle is a fundamental counting technique used to find the number of elements in the union of multiple sets. It states that for any finite sets $A_1, A_2, ..., A_n$, the number of elements in their union can be calculated as follows:

\begin{align*}
    |A_1 \cup A_2 \cup ... \cup A_n| &= \sum_{i=1}^{n} |A_i| \\
    &- \sum_{1 \leq i < j \leq n} |A_i \cap A_j| \\
    &+ \sum_{1 \leq i < j < k \leq n} |A_i \cap A_j \cap A_k| \\
    &- ... + (-1)^{n+1} |A_1 \cap A_2 \cap ... \cap A_n|
\end{align*}

In simpler terms: add the sizes of intersections of an \textbf{odd} number of sets; subtract those of an \textbf{even} number of sets. Equivalently, iterate over all $2^n - 1$ non-empty subsets $S \subseteq \{A_1,\ldots,A_n\}$ and add $(-1)^{|S|+1}|\bigcap_{i \in S} A_i|$.

\textbf{Typical use:} count elements satisfying at least one of $n$ properties (complements: count those satisfying none). Common in problems involving divisibility, forbidden configurations, or surjective functions.